% Ad Atticum 14.18 (ad-atticum-14-18) --- English translation
% Translated by Claude (Anthropic) from the Latin (Perseus).
% Style guide: STYLE.md.

\ciceroLetterOpener{CICERO ATTICO salutem}

\ciceroSection{1} You keep at me now, again and again, that I seem to lift Dolabella's exploit too high into the heavens. As for me --- though I do indeed thoroughly approve of what he did --- still, what drew me to praise him so warmly was your own letters, one and then a second. But Dolabella has now estranged himself wholly from you, and for the very same reason that has made him my own bitterest enemy as well. The shameless man! He owed it on the Kalends of January, and to this day he has not paid; and that after he had freed himself, by Faberius's pen, from an enormous load of debt, and sought aid from Ops (\textit{opem ab Ope}). I am allowed a joke --- you mustn't think I am altogether undone. As it happens, I had sent him a letter on the eighth before the Ides, in the early morning; and on the same day, in the evening, I received yours, here at the Pompeian villa --- swiftly enough, two days out from you. But, as I wrote you that very same day, the letter I sent off to Dolabella has barbs enough on it; and if it produces nothing, I rather think he will not stand up to me face to face.

\ciceroSection{2} The Albian business you have, I take it, settled. As to the Patulcian debt --- which has been, you tell me, \textit{†hung up†} on my account --- it is most welcome, and like everything you do. But I had thought I left Eros behind precisely to clear up that sort of thing; and on his side, not without much fault, the matter has wobbled. I shall see him myself, however.

\ciceroSection{3} About Montanus, as I have written to you so often, do please take the whole affair into your charge. That Servius spoke with you despairingly as he set off does not surprise me in the least, and I concede him nothing on the despair.

\ciceroSection{4} Our Brutus, that singular man --- if he is not to come into the Senate on the Kalends of June, I do not know what he is to do in the Forum. But that he himself knows better. For my part, from what I see being put in train, I judge that the Ides of March did not get us very far. And so I think more and more, day by day, of Greece. For my Brutus, who is meditating exile (so he writes himself), I do not see how I can be of much use. Leonides's letter did not greatly please me. About Herodes I agree with you. I should have liked to read what Saufeius wrote. I was thinking of leaving the Pompeian villa on the sixth before the Ides of May.
